\documentclass[11pt,a4paper]{article}

\usepackage[utf8]{inputenc}
\usepackage[T1]{fontenc}
\usepackage[french]{babel}
\usepackage[a4paper,margin=2.2cm]{geometry}
\usepackage{amsmath,amssymb,amsthm,mathtools}
\usepackage{lmodern}
\usepackage{enumitem}
\usepackage{xcolor}
\usepackage{hyperref}
\usepackage{fancyhdr}
\usepackage{tcolorbox}
\usepackage{array}
\usepackage{longtable}

\hypersetup{
    colorlinks=true,
    linkcolor=black,
    urlcolor=blue
}

\setlength{\parindent}{0pt}
\setlength{\parskip}{0.6em}

\pagestyle{fancy}
\fancyhf{}
\fancyhead[L]{Corrigé détaillé --- DM0 de Physique-Chimie}
\fancyhead[R]{Terminale générale}
\fancyfoot[C]{\thepage}

\newtheorem{remark}{Remarque}
\newcommand{\vect}[1]{\overrightarrow{#1}}

\begin{document}

\begin{center}
    {\LARGE \textbf{Corrigé détaillé du DM de Physique-Chimie}}\\[0.5em]
    {\large \textbf{Avec barème pondéré}}\\[0.5em]
\end{center}

\section*{Barème }

Le sujet est noté sur \textbf{52 points}, avec un barème valorisant davantage les questions longues, ouvertes ou demandant une rédaction approfondie.

\begin{center}
\renewcommand{\arraystretch}{1.3}
\begin{tabular}{|>{\centering\arraybackslash}m{4cm}|>{\centering\arraybackslash}m{3cm}|>{\centering\arraybackslash}m{5cm}|}
\hline
\textbf{Exercice} & \textbf{Barème} & \textbf{Type de compétences valorisées} \\
\hline
Exercice 1 & 9 points & Titrage, raisonnement, protocole \\
\hline
Exercice 2 & 10 points & Analyse de spectre, interprétation \\
\hline
Exercice 3 & 10 points & Modélisation, diffraction, protocole \\
\hline
Exercice 4 & 17 points & Mécanique, démonstrations, discussion ouverte \\
\hline
Exercice 5 & 6 points & Ondes, fréquence, interprétation qualitative \\
\hline
Total & 52 points &  \\
\hline
\end{tabular}
\end{center}

\section*{Exercice 1 --- Titrage et validité d'un protocole \hfill /9}

La réaction support du titrage est :
\[
\mathrm{AH_{(aq)} + HO^-_{(aq)} \rightarrow A^-_{(aq)} + H_2O_{(l)}}.
\]

\subsection*{1. Définir précisément l'équivalence dans un titrage. \hfill /1}

L'équivalence d'un titrage est atteinte lorsque les réactifs ont été introduits dans les \textbf{proportions stœchiométriques} de la réaction support.

Ici, l'acide AH est monoprotique et réagit avec les ions hydroxyde selon un rapport stœchiométrique \(1:1\). À l'équivalence, on a donc :
\[
n(\mathrm{AH})_{\text{initial}} = n(\mathrm{HO^-})_{\text{versés}}.
\]

Autrement dit, la quantité de base versée est exactement celle qu'il faut pour consommer toute la quantité initiale d'acide.

\subsection*{2. Montrer qu'à l'équivalence, \(C_A = \dfrac{C_BV_E}{V_A}\). \hfill /1{,}5}

On note :
\begin{itemize}
    \item \(C_A\) la concentration de la solution acide,
    \item \(V_A\) le volume prélevé de solution acide,
    \item \(C_B\) la concentration de la solution basique,
    \item \(V_E\) le volume de base versé à l'équivalence.
\end{itemize}

La quantité de matière initiale d'acide vaut :
\[
n(\mathrm{AH}) = C_A V_A.
\]

La quantité de matière d'ions hydroxyde versée à l'équivalence vaut :
\[
n(\mathrm{HO^-}) = C_B V_E.
\]

À l'équivalence, les réactifs sont dans les proportions stœchiométriques \(1:1\), donc :
\[
C_A V_A = C_B V_E.
\]

On en déduit :
\[
\boxed{C_A = \frac{C_BV_E}{V_A}}.
\]

\subsection*{3. « Si on double le volume prélevé \(V_A\), alors le volume équivalent double nécessairement. » \hfill /1{,}5}

Cette affirmation est \textbf{vraie}, à condition que :
\begin{itemize}
    \item la solution acide soit la même, donc de concentration \(C_A\) inchangée ;
    \item la solution basique utilisée soit la même, donc de concentration \(C_B\) inchangée.
\end{itemize}

En effet, d'après la relation précédente :
\[
V_E = \frac{C_A V_A}{C_B}.
\]

Si l'on remplace \(V_A\) par \(2V_A\), alors :
\[
V_E' = \frac{C_A(2V_A)}{C_B} = 2V_E.
\]

Donc \textbf{si on double le volume prélevé, alors le volume à l'équivalence double aussi}.

\subsection*{4. Effet des deux choix proposés sur \(V_E\) et sur la précision du titrage. \hfill /2}

\textbf{Proposition de l'élève 1 : diluer la solution acide avant prélèvement}

Si la solution acide est diluée, sa concentration \(C_A\) diminue. Or :
\[
V_E = \frac{C_A V_A}{C_B}.
\]
Donc si \(C_A\) diminue, alors \(V_E\) diminue aussi.

Un volume équivalent plus petit est généralement \textbf{moins favorable à la précision}, car l'incertitude de lecture à la burette représente alors une part plus importante du volume total versé.

\medskip

\textbf{Proposition de l'élève 2 : utiliser une solution basique moins concentrée}

Si \(C_B\) diminue, alors :
\[
V_E = \frac{C_A V_A}{C_B}
\]
augmente.

Un volume équivalent plus grand est souvent \textbf{plus favorable à la précision}, car l'incertitude absolue de lecture sur la burette devient relativement moins importante.

\medskip

\textbf{Conclusion :} pour obtenir un dosage plus précis, il est en général préférable d'utiliser une \textbf{solution titrante moins concentrée} plutôt que de diluer la solution à doser avant prélèvement.

\subsection*{5. Proposer un protocole réaliste pour préparer une dilution au dixième. \hfill /3}

Une dilution par un facteur 10 signifie que la concentration finale est dix fois plus faible que la concentration initiale :
\[
C_f = \frac{C_m}{10}.
\]

On peut procéder ainsi :

\begin{enumerate}[label=\arabic*.]
    \item Prélever \(\,10{,}0\,\mathrm{mL}\,\) de solution mère à l'aide d'une \textbf{pipette jaugée de \(10{,}0\,\mathrm{mL}\)} et d'une propipette.
    \item Introduire ce volume dans une \textbf{fiole jaugée de \(100{,}0\,\mathrm{mL}\)}.
    \item Ajouter un peu d'eau distillée.
    \item Agiter légèrement pour homogénéiser.
    \item Compléter à l'eau distillée jusqu'au \textbf{trait de jauge}.
    \item Boucher la fiole puis homogénéiser par retournements.
\end{enumerate}

\textbf{Précautions expérimentales :}
\begin{itemize}
    \item rincer la pipette avec un peu de solution mère avant prélèvement ;
    \item lire le bas du ménisque à hauteur d'œil ;
    \item ne pas dépasser le trait de jauge ;
    \item homogénéiser soigneusement la solution finale.
\end{itemize}

\textbf{Conclusion :} ce protocole permet bien d'obtenir une dilution au dixième.

\newpage
\section*{Exercice 2 --- Lecture graphique d'un spectre IR et interprétation \hfill /10}

On observe sur le spectre simplifié :
\begin{itemize}
    \item une \textbf{bande large} ;
    \item un \textbf{pic vers \(1700\ \mathrm{cm}^{-1}\)}.
\end{itemize}

Les familles proposées sont :
\begin{itemize}
    \item alcool ;
    \item acide carboxylique ;
    \item ester ;
    \item cétone.
\end{itemize}

\subsection*{1. Relever les deux zones les plus remarquables. \hfill /1}

Les deux zones les plus remarquables sont :
\begin{itemize}
    \item une \textbf{bande large} dans la zone des grands nombres d'onde ;
    \item une \textbf{bande forte vers \(1700\ \mathrm{cm}^{-1}\)}.
\end{itemize}

\subsection*{2. Associer à chacune de ces zones un groupe caractéristique probable. \hfill /1{,}5}

La bande large est caractéristique d'une liaison \(\mathrm{O-H}\).

La bande forte vers \(1700\ \mathrm{cm}^{-1}\) est caractéristique d'un groupe \(\mathrm{C=O}\), appelé \textbf{groupe carbonyle}.

\subsection*{3. Pourquoi la présence d'une bande vers \(1700\ \mathrm{cm}^{-1}\) ne suffit-elle pas à conclure à une cétone ? \hfill /1{,}5}

Une bande vers \(1700\ \mathrm{cm}^{-1}\) indique la présence d'un groupe carbonyle \(\mathrm{C=O}\), mais ce groupe n'est pas spécifique aux cétones.

En effet, un groupe carbonyle est aussi présent dans :
\begin{itemize}
    \item les acides carboxyliques ;
    \item les esters ;
    \item les aldéhydes.
\end{itemize}

Ainsi, \textbf{la seule présence d'une bande vers \(1700\ \mathrm{cm}^{-1}\) ne permet pas d'identifier la famille exacte}. Il faut utiliser d'autres bandes du spectre.

\subsection*{4. Éliminer, en justifiant, les familles incompatibles avec le spectre. \hfill /2}

\textbf{Alcool :} un alcool possède une liaison \(\mathrm{O-H}\), donc peut présenter une bande large, mais il ne possède pas de groupe carbonyle \(\mathrm{C=O}\). Or le spectre présente un pic vers \(1700\ \mathrm{cm}^{-1}\). L'alcool est donc incompatible.

\textbf{Ester :} un ester possède un groupe carbonyle \(\mathrm{C=O}\), mais pas de liaison \(\mathrm{O-H}\) responsable d'une large bande comme celle observée. L'ester est donc incompatible.

\textbf{Cétone :} une cétone possède bien un groupe \(\mathrm{C=O}\), mais pas de liaison \(\mathrm{O-H}\). La présence de la bande large rend donc la cétone incompatible.

\textbf{Acide carboxylique :} un acide carboxylique possède à la fois :
\begin{itemize}
    \item un groupe \(\mathrm{C=O}\) ;
    \item une liaison \(\mathrm{O-H}\) donnant une large bande.
\end{itemize}
Cette famille est donc compatible avec le spectre.

\subsection*{5. Déterminer la famille fonctionnelle la plus probable pour le composé \(X\). \hfill /1}

La famille fonctionnelle la plus probable est celle des \textbf{acides carboxyliques}.

\[
\boxed{\text{Le composé }X\text{ est très probablement un acide carboxylique.}}
\]

\subsection*{6. Pourquoi ce spectre ne permet-il pas forcément d'identifier de manière unique la molécule exacte ? \hfill /1}

Un spectre IR permet surtout d'identifier des \textbf{groupes fonctionnels}. Or plusieurs molécules différentes peuvent posséder les mêmes groupes caractéristiques.

Ainsi :
\begin{itemize}
    \item plusieurs acides carboxyliques auraient un spectre globalement compatible ;
    \item le spectre fourni est simplifié ;
    \item on ne dispose ni de la formule brute, ni de la masse molaire, ni d'un autre spectre complémentaire.
\end{itemize}

Donc ce spectre permet de proposer une \textbf{famille fonctionnelle}, mais pas forcément la molécule exacte.

\subsection*{7. Distinguer expérimentalement un alcool et un acide carboxylique. \hfill /2}

Une stratégie simple consiste à utiliser un \textbf{test acido-basique}.

\textbf{Étape 1 : mesurer le pH}

On mesure le pH des solutions.
\begin{itemize}
    \item Un \textbf{acide carboxylique} donne une solution acide, donc un pH inférieur à \(7\).
    \item Un \textbf{alcool} donne en général une solution de pH voisin de \(7\).
\end{itemize}

\textbf{Étape 2 : ajouter de l'hydrogénocarbonate de sodium}

On ajoute un peu d'hydrogénocarbonate de sodium.
\begin{itemize}
    \item Avec un acide carboxylique, on observe un \textbf{dégagement gazeux} de \(\mathrm{CO_2}\), donc une effervescence.
    \item Avec un alcool, il n'y a pas cette réaction.
\end{itemize}

\textbf{Conclusion :} le liquide qui provoque une effervescence avec l'hydrogénocarbonate et présente un caractère acide est l'acide carboxylique. L'autre liquide est l'alcool.

\section*{Exercice 3 --- Diffraction : établir un modèle puis l'exploiter \hfill /10}

On admet que, pour le premier minimum de diffraction :
\[
\theta \approx \frac{\lambda}{a}.
\]

\subsection*{1. À l'aide du schéma, expliquer pourquoi on peut relier \(L\), \(D\) et \(\theta\). \hfill /1}

La tache centrale s'étend entre les deux premiers minima de diffraction, placés symétriquement par rapport à l'axe central.

Si \(y\) désigne la demi-largeur de la tache centrale, alors :
\[
\tan \theta = \frac{y}{D}.
\]

Comme la largeur totale de la tache centrale vaut \(L = 2y\), on obtient :
\[
L = 2D \tan \theta.
\]

\subsection*{2. Montrer que, pour les petits angles, \(L \approx \dfrac{2\lambda D}{a}\). \hfill /2}

On a :
\[
L = 2D \tan \theta.
\]

Pour les petits angles :
\[
\tan \theta \approx \theta.
\]

Donc :
\[
L \approx 2D\theta.
\]

Or l'énoncé donne :
\[
\theta \approx \frac{\lambda}{a}.
\]

En remplaçant \(\theta\), on obtient :
\[
L \approx 2D \times \frac{\lambda}{a}.
\]

Finalement :
\[
\boxed{L \approx \frac{2\lambda D}{a}}.
\]

\subsection*{3. Sans calcul numérique, indiquer comment varie \(L\). \hfill /1{,}5}

On utilise :
\[
L \approx \frac{2\lambda D}{a}.
\]

\begin{itemize}
    \item Si la largeur \(a\) de la fente augmente, alors \(L\) \textbf{diminue}.
    \item Si la distance \(D\) à l'écran augmente, alors \(L\) \textbf{augmente}.
    \item Si la longueur d'onde \(\lambda\) augmente, alors \(L\) \textbf{augmente}.
\end{itemize}

\subsection*{4. Un élève affirme : « si on divise la largeur de la fente par 2, la figure de diffraction est deux fois moins large ». \hfill /1{,}5}

Cette affirmation est \textbf{fausse}.

En effet :
\[
L \approx \frac{2\lambda D}{a}.
\]

Si on remplace \(a\) par \(\dfrac{a}{2}\), alors :
\[
L' = \frac{2\lambda D}{a/2} = 2 \times \frac{2\lambda D}{a} = 2L.
\]

La tache centrale devient donc \textbf{deux fois plus large}, et non deux fois moins large.

L'argument de l'élève, fondé sur l'idée que ``moins de lumière passe'', ne permet pas de conclure sur la largeur de la figure : cela concerne plutôt l'intensité lumineuse, pas l'extension de la diffraction.

\subsection*{5. Comparer les largeurs des taches centrales. \hfill /1{,}5}

\textbf{Expérience A :} laser rouge et fente de largeur \(a\)
\[
L_A \approx \frac{2\lambda_r D}{a}.
\]

\textbf{Expérience B :} laser vert et fente de largeur \(2a\)
\[
L_B \approx \frac{2\lambda_v D}{2a} = \frac{\lambda_v D}{a}.
\]

On compare :
\[
\frac{L_B}{L_A} = \frac{\lambda_v D/a}{2\lambda_r D/a} = \frac{\lambda_v}{2\lambda_r}.
\]

Or la longueur d'onde du vert est plus petite que celle du rouge, donc :
\[
\lambda_v < \lambda_r \quad \Rightarrow \quad \frac{\lambda_v}{2\lambda_r} < \frac{1}{2}.
\]

Ainsi :
\[
L_B < L_A.
\]

\textbf{Conclusion :} la tache centrale de l'expérience B est plus étroite que celle de l'expérience A.

\subsection*{6. Mesurer expérimentalement la largeur d'un cheveu à l'aide d'un laser. \hfill /2{,}5}

\textbf{Principe physique :}

Un cheveu diffracte la lumière. Au niveau du programme, on peut exploiter une relation du même type que pour une fente :
\[
L \approx \frac{2\lambda D}{a},
\]
où \(a\) désigne ici le diamètre du cheveu.

\textbf{Protocole :}
\begin{enumerate}[label=\arabic*.]
    \item Placer un cheveu verticalement devant un faisceau laser monochromatique.
    \item Placer un écran à une distance \(D\) connue du cheveu.
    \item Observer la figure de diffraction.
    \item Mesurer la largeur \(L\) de la tache centrale.
    \item Utiliser la longueur d'onde \(\lambda\) du laser.
\end{enumerate}

\textbf{Mesures à effectuer :}
\begin{itemize}
    \item la distance \(D\) entre le cheveu et l'écran ;
    \item la largeur \(L\) de la tache centrale ;
    \item la longueur d'onde \(\lambda\) du laser.
\end{itemize}

\textbf{Exploitation :}

À partir de
\[
L \approx \frac{2\lambda D}{a},
\]
on obtient :
\[
\boxed{a \approx \frac{2\lambda D}{L}}.
\]

\textbf{Remarque :} pour améliorer la précision, on peut effectuer plusieurs mesures et faire une moyenne.

\newpage
\section*{Exercice 4 --- Projectile : modèle, trajectoire et discussion des hypothèses \hfill /17}

On étudie le mouvement d'un projectile lancé depuis le point \(A(0,h)\) avec une vitesse initiale de norme \(v_0\), faisant un angle \(\alpha\) avec l'horizontale.  
Dans un premier temps, on néglige les frottements de l'air.

\subsection*{1. Donner les coordonnées du vecteur vitesse initiale. \hfill /1}

La vitesse initiale se décompose en :
\[
\boxed{\vec v_0=
\begin{pmatrix}
v_0\cos\alpha\\
v_0\sin\alpha
\end{pmatrix}}
\]

Ainsi :
\[
v_{0x}=v_0\cos\alpha
\qquad \text{et} \qquad
v_{0z}=v_0\sin\alpha.
\]

\subsection*{2. Montrer que l'accélération du projectile est constante et déterminer ses coordonnées. \hfill /1{,}5}

En négligeant les frottements de l'air, la seule force exercée sur le projectile est son poids :
\[
\vec P = m\vec g.
\]

Dans le repère choisi, avec l'axe vertical vers le haut :
\[
\vec g=
\begin{pmatrix}
0\\
-g
\end{pmatrix}.
\]

La deuxième loi de Newton donne :
\[
m\vec a = \sum \vec F = \vec P = m\vec g.
\]

Donc :
\[
\vec a = \vec g =
\begin{pmatrix}
0\\
-g
\end{pmatrix}.
\]

L'accélération est donc \textbf{constante}.

\subsection*{3. En déduire les expressions de \(x(t)\) et \(z(t)\). \hfill /2}

\textbf{Sur l'axe horizontal :}
\[
a_x=0 \quad \Rightarrow \quad v_x(t)=v_0\cos\alpha.
\]
Comme \(x(0)=0\), on obtient :
\[
\boxed{x(t)=v_0\cos\alpha \, t.}
\]

\textbf{Sur l'axe vertical :}
\[
a_z=-g \quad \Rightarrow \quad v_z(t)=v_0\sin\alpha - gt.
\]
Comme \(z(0)=h\), on obtient :
\[
\boxed{z(t)=h+v_0\sin\alpha \, t - \frac12 gt^2.}
\]

\subsection*{4. Donner l'équation de la trajectoire. Quel est le type de cette trajectoire ? \hfill /2}

On part de :
\[
x(t)=v_0\cos\alpha \, t
\quad \Rightarrow \quad
t=\frac{x}{v_0\cos\alpha}.
\]

On remplace \(t\) dans l'expression de \(z(t)\) :
\[
z(x)=h+v_0\sin\alpha \cdot \frac{x}{v_0\cos\alpha}
-\frac12 g\left(\frac{x}{v_0\cos\alpha}\right)^2.
\]

Donc :
\[
z(x)=h+x\tan\alpha-\frac{g}{2v_0^2\cos^2\alpha}x^2.
\]

Ainsi :
\[
\boxed{z(x)=h+x\tan\alpha-\frac{g}{2v_0^2\cos^2\alpha}x^2.}
\]

Cette expression est un polynôme du second degré en \(x\), donc la trajectoire est une \textbf{parabole}.

\subsection*{5. Expliquer pourquoi le mouvement horizontal est uniforme alors que le mouvement vertical ne l'est pas. \hfill /1{,}5}

Le mouvement horizontal est uniforme parce que :
\begin{itemize}
    \item aucune force n'agit horizontalement ;
    \item donc \(a_x=0\) ;
    \item la vitesse horizontale \(v_x\) est constante.
\end{itemize}

Le mouvement vertical n'est pas uniforme parce que :
\begin{itemize}
    \item le poids agit verticalement ;
    \item donc \(a_z=-g\neq 0\) ;
    \item la vitesse verticale varie au cours du temps.
\end{itemize}

\subsection*{6. « Au point le plus haut, la vitesse du projectile est nulle. » \hfill /1{,}5}

Cette affirmation est \textbf{fausse}.

Au point le plus haut, seule la composante verticale de la vitesse s'annule :
\[
v_z=0.
\]

Mais la composante horizontale reste constante :
\[
v_x=v_0\cos\alpha.
\]

Donc le vecteur vitesse au sommet est :
\[
\vec v_{\text{sommet}}=
\begin{pmatrix}
v_0\cos\alpha\\
0
\end{pmatrix}.
\]

La vitesse n'est donc pas nulle, sauf dans le cas particulier d'un lancer strictement vertical.

\subsection*{7. Déterminer littéralement l'instant auquel l'altitude est maximale. \hfill /1{,}5}

L'altitude est maximale lorsque la vitesse verticale s'annule :
\[
v_z(t)=v_0\sin\alpha - gt = 0.
\]

On obtient :
\[
\boxed{t_{\max}=\frac{v_0\sin\alpha}{g}.}
\]

\subsection*{8. En déduire littéralement la hauteur maximale atteinte. \hfill /2}

On remplace \(t_{\max}\) dans l'expression de \(z(t)\) :
\[
z_{\max}=h+v_0\sin\alpha \cdot \frac{v_0\sin\alpha}{g}
-\frac12 g\left(\frac{v_0\sin\alpha}{g}\right)^2.
\]

D'où :
\[
z_{\max}=h+\frac{v_0^2\sin^2\alpha}{g}
-\frac12\frac{v_0^2\sin^2\alpha}{g}.
\]

Finalement :
\[
\boxed{z_{\max}=h+\frac{v_0^2\sin^2\alpha}{2g}.}
\]

\subsection*{9. Méthode générale pour savoir si le projectile franchit un obstacle. \hfill /1{,}5}

Supposons qu'un obstacle soit situé à l'abscisse \(x=x_{\text{obs}}\) et ait une hauteur \(H\).

La méthode consiste à :
\begin{enumerate}[label=\arabic*.]
    \item établir l'équation de la trajectoire
    \[
    z(x)=h+x\tan\alpha-\frac{g}{2v_0^2\cos^2\alpha}x^2 ;
    \]
    \item calculer la hauteur du projectile au niveau de l'obstacle :
    \[
    z(x_{\text{obs}}) ;
    \]
    \item comparer cette valeur à la hauteur \(H\).
\end{enumerate}

\begin{itemize}
    \item Si \(z(x_{\text{obs}})>H\), le projectile franchit l'obstacle.
    \item Si \(z(x_{\text{obs}})=H\), il le touche exactement.
    \item Si \(z(x_{\text{obs}})<H\), il ne le franchit pas.
\end{itemize}

Il faut en outre vérifier que l'obstacle se situe bien dans une zone effectivement atteinte pendant le vol.

\subsection*{10. Discussion argumentée : dans quelles situations le modèle sans frottements est-il acceptable ? \hfill /3}

\textbf{Sens physique de l'hypothèse}

Négliger les frottements de l'air signifie que l'on suppose que l'air n'exerce pas de force significative sur le projectile. La seule force prise en compte est alors le poids.

\medskip

\textbf{Situations où cette hypothèse peut paraître raisonnable}

Cette hypothèse peut être acceptable :
\begin{itemize}
    \item pour des objets assez massifs ;
    \item pour des vitesses modérées ;
    \item pour des trajectoires courtes ;
    \item lorsque l'on cherche seulement une première modélisation.
\end{itemize}

Exemples :
\begin{itemize}
    \item lancer d'une balle dense sur une courte distance ;
    \item expérience de laboratoire simplifiée ;
    \item exercice théorique destiné à comprendre les lois du mouvement.
\end{itemize}

\medskip

\textbf{Situations où cette hypothèse devient discutable}

Elle devient discutable lorsque :
\begin{itemize}
    \item la vitesse est élevée ;
    \item la durée du vol est longue ;
    \item l'objet est léger ;
    \item la surface de contact avec l'air est importante ;
    \item il y a du vent.
\end{itemize}

Exemples :
\begin{itemize}
    \item balle de ping-pong ;
    \item plume ;
    \item feuille de papier ;
    \item ballon ralenti par l'air ;
    \item jet d'eau fin.
\end{itemize}

\medskip

\textbf{Grandeurs modifiées par les frottements}

La présence de frottements peut modifier :
\begin{itemize}
    \item la portée ;
    \item la durée du vol ;
    \item la hauteur maximale ;
    \item la valeur de la vitesse au cours du mouvement ;
    \item la forme de la trajectoire, qui n'est plus exactement parabolique.
\end{itemize}

\medskip

\textbf{Conclusion}

Le modèle sans frottements est \textbf{utile} car il permet une description simple, claire et souvent suffisante pour une première étude. Mais il possède des \textbf{limites} : dans de nombreuses situations réelles, les frottements de l'air jouent un rôle non négligeable, et un modèle plus élaboré devient alors nécessaire.

\newpage
\section*{Exercice 5 --- Pourquoi une sirène paraît-elle différente quand elle s'éloigne ? \hfill /6}

On étudie qualitativement le son émis par une sirène de véhicule en mouvement.  
Un observateur immobile perçoit un son \textbf{plus aigu} lorsque le véhicule s'approche, puis \textbf{plus grave} lorsqu'il s'éloigne.

\subsection*{Corrigé détaillé}

Le phénomène décrit correspond à l'\textbf{effet Doppler}.

\medskip

\textbf{Idée générale :} lorsqu'une source sonore est en mouvement par rapport à un observateur, la fréquence du son \textbf{perçue} par l'observateur n'est pas la même que la fréquence \textbf{émise} par la source.

\medskip

\textbf{Quand la sirène s'approche :}

La source se rapproche de l'observateur. Les fronts d'onde successifs arrivent alors plus rapprochés dans l'espace au niveau de l'observateur.  
Autrement dit, l'intervalle de temps entre l'arrivée de deux signaux identiques successifs devient plus petit.

La période perçue diminue donc :
\[
T_{\text{perçue}} < T_{\text{émise}}.
\]

Or la fréquence est l'inverse de la période :
\[
f=\frac{1}{T}.
\]

Donc la fréquence perçue augmente :
\[
f_{\text{perçue}} > f_{\text{émise}}.
\]

Comme un son plus aigu correspond à une fréquence plus élevée, l'observateur entend donc une sirène \textbf{plus aiguë} lorsque le véhicule s'approche.

\medskip

\textbf{Quand la sirène s'éloigne :}

La source s'éloigne de l'observateur. Les fronts d'onde successifs arrivent alors plus espacés.

L'intervalle de temps entre deux motifs identiques successifs augmente. La période perçue devient donc plus grande :
\[
T_{\text{perçue}} > T_{\text{émise}}.
\]

Par conséquent, la fréquence perçue diminue :
\[
f_{\text{perçue}} < f_{\text{émise}}.
\]

Le son paraît alors \textbf{plus grave}.

\medskip

\textbf{Interprétation en termes de signal périodique :}

Un son musical ou une sirène peuvent être modélisés par un signal périodique.  
Lorsque la source est immobile par rapport à l'observateur, la fréquence perçue est la même que la fréquence émise.

En revanche, si la source se déplace :
\begin{itemize}
    \item vers l'observateur, le signal reçu semble plus \og resserré \fg{} dans le temps ;
    \item en s'éloignant, le signal reçu semble plus \og étiré \fg{} dans le temps.
\end{itemize}

Cela se traduit par une modification de la période et donc de la fréquence perçue.

\medskip

\textbf{Conclusion :}

La sirène paraît plus aiguë quand elle s'approche et plus grave quand elle s'éloigne parce que le mouvement de la source modifie la fréquence perçue par l'observateur.  
Ce phénomène est appelé \textbf{effet Doppler}.

\subsection*{Éléments attendus dans une bonne réponse}

Une réponse bien rédigée devait faire apparaître :
\begin{itemize}
    \item la notion de \textbf{signal périodique} ;
    \item le lien entre \textbf{période} et \textbf{fréquence} :
    \[
    f=\frac{1}{T} ;
    \]
    \item l'idée que lorsque la source s'approche, les ondes sont reçues plus rapprochées ;
    \item l'idée que lorsque la source s'éloigne, les ondes sont reçues plus espacées ;
    \item la conclusion : \textbf{son plus aigu à l'approche, plus grave à l'éloignement}.
\end{itemize}

\section*{Barème détaillé question par question}

\begin{longtable}{|p{9cm}|c|}
\hline
\textbf{Question} & \textbf{Points} \\
\hline
\endfirsthead
\hline
\textbf{Question} & \textbf{Points} \\
\hline
\endhead

Exercice 1, question 1 & 1 \\
\hline
Exercice 1, question 2 & 1{,}5 \\
\hline
Exercice 1, question 3 & 1{,}5 \\
\hline
Exercice 1, question 4 & 2 \\
\hline
Exercice 1, question 5 & 3 \\
\hline

Exercice 2, question 1 & 1 \\
\hline
Exercice 2, question 2 & 1{,}5 \\
\hline
Exercice 2, question 3 & 1{,}5 \\
\hline
Exercice 2, question 4 & 2 \\
\hline
Exercice 2, question 5 & 1 \\
\hline
Exercice 2, question 6 & 1 \\
\hline
Exercice 2, question 7 & 2 \\
\hline

Exercice 3, question 1 & 1 \\
\hline
Exercice 3, question 2 & 2 \\
\hline
Exercice 3, question 3 & 1{,}5 \\
\hline
Exercice 3, question 4 & 1{,}5 \\
\hline
Exercice 3, question 5 & 1{,}5 \\
\hline
Exercice 3, question 6 & 2{,}5 \\
\hline

Exercice 4, question 1 & 1 \\
\hline
Exercice 4, question 2 & 1{,}5 \\
\hline
Exercice 4, question 3 & 2 \\
\hline
Exercice 4, question 4 & 2 \\
\hline
Exercice 4, question 5 & 1{,}5 \\
\hline
Exercice 4, question 6 & 1{,}5 \\
\hline
Exercice 4, question 7 & 1{,}5 \\
\hline
Exercice 4, question 8 & 2 \\
\hline
Exercice 4, question 9 & 1{,}5 \\
\hline
Exercice 4, question 10 & 3 \\
\hline

Exercice 5 --- explication qualitative complète de l'effet Doppler & 6 \\
\hline

\textbf{Total} & \textbf{52} \\
\hline
\end{longtable}

\end{document}